\subsection{REST Error Codes}

All REST API endpoints return errors as JSON objects in the format \texttt{\{"error": "message"\}}. The following table summarizes the HTTP status codes used across the API and their meaning.

\begin{longtable}{|p{.15\columnwidth}|p{.2\columnwidth} |p{.6\columnwidth}|}
\hline
\textbf{Error Code} & \textbf{HTTP Status} & \textbf{Description} \\\hline

400 & Bad Request & Malformed JSON body, missing required fields, invalid UUID format, invalid amount or date format, empty request body, oversized request body (>8KB), invalid scenario value, or blank required parameters. \\\hline

401 & Unauthorized & Authentication required. Returned by \texttt{AuthenticationFilter} for unauthenticated \texttt{/api/*} requests, or by the servlet-level defensive check if the session has expired. \\\hline

403 & Forbidden & The authenticated user does not own the requested resource (e.g., attempting to modify another user's wallet or transaction). \\\hline

404 & Not Found & The requested resource does not exist (e.g., no forecast computed yet for the given scenario). \\\hline

500 & Internal Server Error & An unexpected server error occurred. Internal details are never exposed; the response contains only a generic error message. \\\hline

\caption{HTTP status codes used by the REST API.}
\label{tab:restErrorCodes}
\end{longtable}
